% Run from the repository root with LuaLaTeX. The same file supports report classes.
\documentclass[fontset=fandol]{hithesis}
\begin{document}
\directlua{dofile('tests/font-baseline.lua')}
\setbox0=\hbox{\songti 国永哈中，ABC $x+1$}
\directlua{baseline_check_box(0,'songti',true)}
\setbox0=\hbox{\heiti 国永哈中，ABC}
\directlua{baseline_check_box(0,'heiti',true)}
\setbox0=\hbox{\kaishu 国永哈中，ABC}
\directlua{baseline_check_box(0,'kaishu',true)}
\setbox0=\hbox{\fangsong 国永哈中，ABC}
\directlua{baseline_check_box(0,'fangsong',true)}
\setbox0=\hbox{\songti\bfseries 国永哈中，ABC}
\directlua{baseline_check_box(0,'songti',true)}
\setbox0=\hbox{\songti\itshape 国永哈中，ABC}
\directlua{baseline_check_box(0,'songti',true)}
\setbox0=\hbox{\songti\fontsize{18bp}{24bp}\selectfont 国永哈中，ABC}
\directlua{baseline_check_box(0,'songti',true)}
\setbox0=\hbox{\ttfamily 国永哈中，ABC}
\directlua{baseline_check_box(0,'fangsong',true)}
\setbox0=\hbox{\songti\emph{国永哈中，ABC}}
\directlua{baseline_check_box(0,'kaishu',true)}
\begingroup
\hitsetup{font-baseline-correction=false}
\setbox0=\hbox{\songti 国永哈中，ABC}
\directlua{baseline_check_box(0,'songti',false)}
\endgroup
\setbox0=\hbox{\songti 国永哈中，ABC}
\directlua{baseline_check_box(0,'songti',true)}
\ExplSyntaxOn
\group_begin:
  \prop_put:Nnn\g__hit_font_roles_prop{songti}{FandolKai}
  \__hit_font_cjk_declare:nn{baseline-shared}{songti}
\group_end:
\ExplSyntaxOff
\setbox0=\hbox{\CJKfamily{baseline-shared}国永哈中，ABC}
\directlua{baseline_check_box(0,'songti',true)}
\setbox0=\hbox{\kaishu 国永哈中，ABC}
\directlua{baseline_check_box(0,'kaishu',true)}
\begingroup
\fontspec{texgyretermes-regular.otf}
\directlua{assert(docxlike.run.baseline_register(font.current(),'no-CJK-anchor',0.3)==0)}
\endgroup
\IfFontExistsTF{SimSun}{\begingroup\fontspec{SimSun}
  \directlua{local i=docxlike.run.baseline_register(font.current(),'Windows-SimSun',0.361328125); assert(i>0 and docxlike.run.baseline_records[i].delta==0)}
\endgroup}{}
\def\baselineparagraph{中文字体基线校正只移动字形，不改变断行位置。
  English text and $a+b=c$ remain unchanged.
  \texttt{print("中文代码")}。\textbf{粗体文字}和\emph{强调文字}。\par}
\hitsetup{font-baseline-correction=false}
\setbox0=\vbox{\hsize=180bp\baselineparagraph\baselineparagraph}
\hitsetup{font-baseline-correction=auto}
\setbox2=\vbox{\hsize=180bp\baselineparagraph\baselineparagraph}
\directlua{baseline_check_layout(0,2)}
\noindent\copy0\quad\copy2
\directlua{baseline_finish()}
\end{document}
